\section{Conclusion}
% PhyPrompt employs a concise two-stage pipeline: first, supervised fine-tuning on a physics-focused Chain-of-Thought corpus, then reinforcement learning with a dynamic reward curriculum that gradually shifts emphasis from semantic fidelity to commonsense physics. This process automatically refines text-to-video prompts for substantially improved physical realism achieving a 40.8\% combined semantic-and-physical success rate, while generalizing zero-shot across diverse generators.
We present PhyPrompt, a reinforcement learning-based framework that automatically enhances text-to-video prompts for physically plausible generation. Through a two-stage training approach combining physics-focused Chain-of-Thought supervision with dynamic reward-based reinforcement learning, PhyPrompt achieves substantial improvements in physical realism while maintaining semantic fidelity.
Our key empirical finding demonstrates that semantic adherence and physical commonsense are not inherently competing objectives when optimized through an appropriate curriculum. PhyPrompt-7B achieves 47.8\% semantic adherence and 66.8\% physical commonsense on the VideoPhy2 benchmark, yielding a 40.8\% joint success rate. Critically, this represents an 11 percentage point improvement in physical commonsense (from 55.8\% to 66.8\%) without sacrificing semantic quality. In fact, semantic adherence simultaneously improves by 4.4 percentage points (from 43.4\% to 47.8\%), resulting in an 8.6 percentage point gain in joint performance over baseline prompts.
This synergistic improvement stems from our dynamic curriculum mechanism, which discovers prompt-space regions unreachable through single-objective optimization. By establishing semantic scaffolding early in training before refining with physical constraints, the curriculum enables both objectives to mutually reinforce rather than trade off against each other. Our ablation studies confirm this effect: the dynamic curriculum outperforms not only static balanced weighting but also specialized single-objective training on each objective's own metric. Beyond performance gains, PhyPrompt demonstrates genuine zero-shot transfer across architecturally diverse generators (Lavie, VideoCrafter2, CogVideoX), achieving consistent improvements without per-generator fine-tuning. This model-agnostic design makes PhyPrompt a practical, deployable solution for physics-aware video generation.
Our results carry broader implications for multi-objective optimization in generative AI. The success of domain-specialized training with direct task feedback over raw parameter scaling suggests that careful curriculum design and targeted optimization can be more effective than simply increasing model size. As text-to-video systems continue advancing, prompt refinement approaches like PhyPrompt offer a parameter-efficient path toward physically grounded video synthesis suitable for robotics, scientific visualization, and educational applications where physical realism is essential.

% \subsection{Broader Impacts}
% \label{sec:impact}
% PhyPrompt offers clear positive benefits by enabling more physically plausible video synthesis, which can enhance applications in robotics, scientific visualization, virtual training, and education where realistic physics are crucial for downstream tasks. By improving semantic‐and‐physical fidelity, our met
